\section{Renormalization and Matching for PDFs}
\label{sec:renorm}
%-----------------------------------------------------------------------------
In this section, we consider the LaMET application to calculating
the simplest collinear PDFs, which have been most extensively studied in the literature so far.
Although universality allows one to extract the collinear PDFs from the matrix elements of a wide class of operators evaluated at large momentum, we will focus on physical observables closely resembling the collinear PDFs, i.e.,
the quark and gluon momentum distributions or the {\it quasi-PDFs}. We also discuss
the coordinate-space factorization approach in which
the pseudo-PDF and current-current correlators have been studied.

We mainly review the technical progress made in renormalization and matching using the quasi-PDFs.
The matching can be done in principle at the bare matrix elements level,
since the factorization formula like \eq{fact_0} is valid for both bare and
renormalized momentum distributions. All the UV divergences in the bare quasi-PDF can be
factorized into the matching coefficient $C$, and the latter automatically
renormalizes the bare lattice matrix elements, so the continuum limit can be taken afterwards.
However,  such a matching coefficient then has to be calculated in lattice perturbation theory, which is computationally challenging and converges slowly~\cite{Lepage:1992xa}. More importantly, the quasi-PDF contains linear power divergence under UV cutoff regularization due to the Wilson line self-energy~\cite{Ji:2013dva,Xiong:2013bka}, which makes it impossible to take the continuum limit with fixed-order calculations in lattice perturbation theory. Though the latter problem can be improved by resumming the linear and possibly logarithmic divergences, it is usually preferred to nonperturbatively renormalize the quasi-PDFs on the lattice, after which a continuum limit can be taken and a perturbative matching can be done in the continuum theory. To this end, a thorough understanding of the renormalizability of Wilson-line operators defining the quasi-PDFs is required.
In addition to renormalization, the applications of LaMET rely on the validity of the large-momentum factorization formula~\eq{fact_0}, which can be proven in perturbation theory to all orders by showing that the collinear divergences
are the same in the momentum distributions and light-cone PDFs.

We begin in \sec{renorm-nl} with the proof of multiplicative renormalizability of the Wilson-line operators that define the quasi-PDFs. We first work in the continuum theory with $\MS$ scheme, and then generalize the conclusion to lattice theory.
Next, in \sec{factorization} we outline the factorization theorem for momentum distributions to all orders
in perturbation theory, and state the form of convolution between the matching coefficient and the PDF. In \sec{renorm-coord} we show that the factorization theorem has an equivalent form in coordinate space, which can be used as an alternate route to extract PDFs from lattice matrix elements. Finally, we discuss the nonperturbative renormalization of quasi-PDFs on the lattice and their matching to the $\MS$ PDF in \sec{renorm-npr}.

%------------------------------------------------------------------------------
\subsection{Renormalization of Nonlocal Wilson-Line Operators}
\label{sec:renorm-nl}
%------------------------------------------------------------------------------
The momentum distributions of the proton are defined from
equal-time nonlocal Wilson line operators of the form in Eq.~(\ref{eq:momdis1}).
In this subsection, we review the renormalization of these spacelike nonlocal operators (the renormalization of lightlike nonlocal operators defining the PDFs can be found in~\cite{Collins:1981uw,Collins:2011zzd}). We first discuss their renormalization in DR using an auxiliary field approach, followed by the discussion on similar
gluon operators. We then consider power divergences
in the momentum cutoff type of UV regularization. The conclusion is that they are
all multiplicatively renormalizable with a finite number of mixings with other operators.

\subsubsection{Renormalization of nonlocal quark operators}

We are interested in operators of the following kind,
\begin{equation}
	O_\Gamma(z)=\bar{\psi}\big({z\over2}\big)\Gamma W\big({z\over2},-{z\over2}\big)\psi\big(\!-\!{z\over2}\big) \ .
\end{equation}
Since the Wilson line $W(z_1,z_2)$ is a path-ordered integral of gauge fields,
it is not obvious that such operators are multiplicatively renormalizable. The renormalization of non-lightlike Wilson loops and Wilson lines has been studied in early literature~\cite{Dotsenko:1979wb,Craigie:1980qs}, and the all-order proof of their multiplicative renormalizability was first made using diagrammatic methods in~\cite{Dotsenko:1979wb,Craigie:1980qs} and then the functional formalism of gauge theories in~\cite{Dorn:1986dt}.
The same conclusion was conjectured to hold also for the quark bilinear operator $O_\Gamma(z)$, whose renormalization takes the following form~\cite{Musch:2010ka,Ishikawa:2016znu,Chen:2016fxx},
\begin{align} \label{eq:schemeren}
	O^B_\Gamma(z,\Lambda)=Z_{\psi,z}(\Lambda,\mu)e^{\delta m(\Lambda)|z|} O^R_\Gamma(z,\mu)\,,
\end{align}
where ``$B$'' and ``$R$'' stand for bare and renormalized operators respectively, and all the fields and couplings in $O^B_\Gamma(z,\Lambda)$ are bare ones which depend on the UV cutoff $\Lambda$. $\delta m(\Lambda)$ is the ``mass correction'' of the Wilson line, which includes all the linear power divergences of its self-energy. $Z_{\psi,z}(\Lambda,\mu)$ includes all the logarithmic divergences from wavefunction and vertex renormalizations.

An early two-loop study of the quasi-PDF in the $\MS$ scheme indeed indicated the multiplicative renormalizability of $O_\Gamma(z)$~\cite{Ji:2015jwa}. The first rigorous proof of \eq{schemeren} was given in the auxiliary ``heavy quark'' field formalism~\cite{Ji:2017oey,Green:2017xeu} which was used to prove the renormalizability of Wilson lines~\cite{Samuel:1978iy,Gervais:1979fv,Arefeva:1980zd,Dorn:1986dt}. This formalism is defined by extending the QCD Lagrangian to include the auxiliary ``heavy quark'' fields $Q,\bar{Q}$ and their gauge interaction,
\begin{align} \label{eq:hqet}
	{\cal L} = {\cal L}_{\rm QCD} + \bar{Q}_0in_z\cdot D_0 Q_0\,,
\end{align}
where the subscript ``0'' denotes bare quantities. $n_z^\mu=(0,0,0,1)$ is the direction vector of the spacelike Wilson line $W(z,0)$, $D_{0}^\mu = \partial^\mu + ig_0A_{0}^\mu$, and $Q_0$ is a color-triplet scalar Grassmann field in the fundamental representation of SU(3). Note that if we replace $n_z^\mu$ with the timelike vector $n_t^\mu=(1,0,0,0)$, then \eq{hqet} yields the leading order HQET Lagrangian.

In the theory defined by \eq{hqet}, the Wilson line can be expressed as the connected two-point function of the ``heavy-quark'' fields,
\begin{align}
	\langle Q_0(\xi)\bar{Q}_0(\eta) \rangle_Q = S_0^Q(\xi,\eta)\,,
\end{align}
where $\xi$ and $\eta$ are space-time coordinates, and
$\langle...\rangle_Q$ stands for integrating out the auxiliary fields. The above equation is valid up to the determinant of $in_z\cdot D_0$, which is a constant and can be absorbed into the normalization of the generating functional~\cite{Mannel:1991mc}. The Green's function $S_0^Q(\xi,\eta)$ satisfies
\begin{align}
	in_z\cdot D_0(\xi)\ S_0^Q(\xi,\eta) = \delta^{(4)}(\xi-\eta)\,,
\end{align}
which has the solution
\begin{align}
	S_0^Q(\xi,\eta)\! =\! W(\xi^3,\eta^3)\theta(\xi^3-\eta^3)\delta(\xi^0-\eta^0)\delta^{(2)}(\vec{\xi}_\perp -\vec{\eta}_\perp)
\end{align}
with a proper choice of boundary condition.
In this way, the Wilson-line operator $O^B_\Gamma(z)$ can be replaced by the product of two local composite operators averaged over all the ``heavy-quark'' field configurations~\cite{Dorn:1986dt},
\begin{align} \label{eq:product}
	O^B_{\Gamma}(z)& = \int d^4\xi\ \delta(\xi^3-z\big) \nn\\
	&\quad\times \langle \bar{\psi}_0\big({\xi\over2}\big) Q_0\big({\xi\over2}\big)\Gamma\bar{Q}_0\big(\!-\!{\xi\over2}\big)\psi_0\big(\!-\!{\xi\over2}\big)\rangle_Q\,,
\end{align}
where the UV regulator is suppressed.

Consequently, the renormalization of $O^B_{\Gamma}(z)$ is reduced to that of the two local ``heavy-to-light'' currents
\begin{align}
	J^B = \bar{Q}_0 \psi_0\,.
\end{align}
The renormalizability of this auxiliary field theory has been proven using the standard functional techniques for gauge theories~\cite{Dorn:1986dt}. After fixing the covariant gauge and introducing the ghost fields, the theory including the auxiliary ``heavy-quark'' has a residual BRST symmetry, from which one can derive the Ward-Takahashi identities to show that all the UV divergences of the Green's functions can be subtracted with a finite number of local counterterms. In analogy, the same method has also been used to prove the all-order renormalization of HQET in perturbation theory~\cite{Bagan:1993zv}.

According to~\cite{Dorn:1986dt}, the ``heavy-quark'' Lagrangian can be renormalized in a covariant gauge as
\begin{align}\label{eq:auxl}
	{\cal L} &= {\cal L}_{\rm QCD}[g_0,\psi_0,A_0,c_0] + \bar{Q}_0in_z\cdot D_0 Q_0\nn\\
	&={\cal L}_{\rm QCD}[g,\psi,A,c] +  {\cal L}_{\rm c.t.}[g,\psi,A,c] \nn\\
	& +  Z_Q\bar{Q}\left(in_z\cdot \partial - i\delta m\right) Q - g Z_{1}^{QQg}\bar{Q}n_z\cdot A_at^aQ\,,
\end{align}
where ${\cal L}_{\rm c.t.}[g,\psi,A,c]$ are the QCD counterterms, and the bare fields and coupling are related to the renormalized ones through
\begin{align}
	\psi_0=Z_\psi^{1\over2}\psi, \ A_0 = Z_A^{1\over2} A,\ Q_0 = Z_Q^{1\over2} Q,\ g_0=Z_gg\,.
\end{align}
The heavy-quark-gluon vertex renormalization constant $Z_1^{QQg}$ is related to $Z_g$ through the Slavnov-Taylor identities of the auxiliary field theory~\cite{Dorn:1986dt},
\beq
Z_g= Z_1^{QQg} Z_A^{-{1\over2}}Z_Q^{-1}	\,.
\eeq
The $i\delta m$ can be regarded as the mass correction of the ``heavy quark'' except that it is imaginary.
For Dirac fermions, the mass correction is logarithmically divergent and proportional to the bare mass, as a result of chiral symmetry; for HQET, the mass correction of the heavy quark is proportional to the UV cutoff, i.e. linearly divergent, which is also expected for the auxiliary field here. Since the proof of renormalizability for this auxiliary field theory is carried out in the $\MS$ scheme with DR ($d=4-2\epsilon$), all power divergences vanish, so does $\delta m$.
Nevertheless, $\delta m$ may include ${\cal O}(\Lambda_{\text{QCD}})$ contributions due to the renormalon ambiguities which are known to exist in HQET~\cite{Bigi:1994em,Beneke:1994sw}.

Since the auxiliary field theory is renormalizable, the renormalization of the operator product in \eq{product} amounts to the renormalizations of the two ``heavy-to-light'' currents. Using the standard techniques in quantum field theory~\cite{Collins:1984xc}, one can show recursively that the overall UV divergence of the insertion of $J^B$ into Green's functions is absorbed into a renormalization factor $Z_J$ to all orders in perturbation theory,
\begin{align}
	J^B&=Z_{J}J^R = Z_{\psi}^{1/2}Z_{Q}^{1/2}Z_{V}\ J^R\,,
\end{align}
where $Z_{V}$ is the vertex renormalization constant of the ``heavy-to-light'' current.
The renormalization of heavy-to-light currents in HQET has been calculated up to three-loop order in perturbative QCD~\cite{Shifman:1986sm,Politzer:1988wp,Ji:1991pr,Chetyrkin:2003vi,Broadhurst:1991fz}. More recently, it has been argued that the anomalous dimension of the ``heavy-to-light'' current is identical to that in HQET to all orders~\cite{Braun:2020ymy}, which is also the case for the ``heavy-to-gluon'' current that will be discussed below, so the renormalization factors for the spacelike and timelike Wilson line operators should be exactly the same.

Using the above results, we can show that
\begin{align} \label{eq:renproof}
	O^B_\Gamma(z)\!&= Z_J^2\!\int\! d^4\xi\ \delta(\xi^3\!-\!z)\big\langle \bar{J}^R\big({\xi\over2}\big) \Gamma  J^R\big(\!-\!{\xi\over2}\big)\big\rangle_Q\nn\\
	&= Z_J^2 e^{\delta m|z|} O^R_\Gamma(z)\,,
\end{align}
where $\delta m$ arises from the determinant of $(in_z\cdot \partial - \delta m)$ in \eq{auxl}. In this way, we identify that $Z_{\psi,z} =  Z_J^2 $ in \eq{schemeren} which is independent of $\Gamma$. At one-loop order~\cite{Stefanis:1983ke},
\beq \label{eq:zpsiz}
Z_{\psi,z} = 1+ \frac{\alpha_sC_F}{4\pi}{3\over\epsilon_{\mbox{\tiny UV}}}\,,
\eeq
where the UV regulator $\epsilon_{\mbox{\tiny UV}}$ is to be distinguished from the IR regulator $\epsilon_{\mbox{\tiny IR}}$ in DR.

The multiplicative renormalizability of $ O^B_\Gamma(z)$ has also been proven with a recursive analysis of all-order Feynman diagrams~\cite{Ishikawa:2017faj}. In addition to \eq{schemeren}, it was found that $O^B_\Gamma(z)$ does not mix with gluons or quarks of other flavors. This can also be easily understood within the auxiliary field formalism, as the flavor-changing ``heavy-to-light'' current does not mix with other operators~\cite{Green:2020xco}.

Finally, under lattice regularization we can still use the above techniques to prove \eq{renproof}, where the mass correction $\delta m$ is now nonvanishing and equal to the lattice UV cutoff $1/a$ multiplied by a perturbative series in the coupling constant $\alpha_s$.

\subsubsection{Renormalization of nonlocal gluon operators}

Using the same ``heavy-quark'' auxiliary field formalism, it has also been proven that the Wilson-line operators for the gluon quasi-PDF are multiplicatively renormalizable~\cite{Zhang:2018diq}, which is echoed by the diagrammatical proof in~\cite{Li:2018tpe}.

According to LaMET, the gluon quasi-PDF can be defined as~\cite{Ji:2013dva}
\begin{align} \label{eq:quasigluon}
	\tilde g(x,P^z) = N\int {d\lambda \over 4\pi x (P^z)^2}e^{i\lambda x} \langle P| O_g(z)|P\rangle\,,
\end{align}
where $N$ is a normalization factor, and
\begin{align}
	O^B_g(z) = g_{\perp,\mu\nu} F^{n_1\mu}_{0,a}\big({z\over2}\big) W^{ab}\big({z\over2},-{z\over2}\big) F^{n_2\nu}_{0,b}\big(\!-\!{z\over2}\big)
\end{align}
with $F^{n\mu}_{0,a}=n_\rho F^{\rho\mu}_{0,a}$ and $n_1^\mu$, $n_2^\mu$ being either $n_z^\mu$ or $n_t^\mu$. $a,b$ are color indices in the adjoint representation. The transverse metric tensor
\begin{align}
	g_\perp^{\mu \nu} = g^{\mu \nu}  - n_t^\mu n_t^\nu / n_t^2 - n_z^\mu n_z^\nu /n_z^2\,,
\end{align}
and $N = (n_z\cdot P/n_t\cdot P)^{(n_1+n_2)\cdot n_t}$.
For lattice implementation, $O^B_g(z)$ can also be defined as~\cite{Dorn:1986dt,Zhang:2018diq}
\begin{align}
	O^B_g(z)\!=\! 2g_{\perp}^{\mu\nu}\mbox{tr}\Big[F^{n_1}_{0,\mu}\big({z\over2}\big)W\!\big({z\over2},\!-\!{z\over2}\big)F^{n_2}_{0,\nu}\big(\!-\!{z\over2}\big)W\!\big(\!-\!{z\over2},\!{z\over2}\big)\Big]\,,
\end{align}
where $F^{\mu\nu}=F^{\mu\nu}_at^a$ and $W$ are in the fundamental representation.
Similar to \eq{product}, we can express $O^B_g(z)$ as a product of two local composite operators,
\begin{align} \label{eq:gluonprod}
	\tilde{O}_g^B(z)&=\int d^4\xi\ \delta(\xi^3\!-\!z)\\
	&\times g_{\perp,\mu\nu}\big\langle F^{n_1\mu}_{0,a}\big({\xi\over2}\big)Q^a_0\big({\xi\over2}\big)\bar{Q}^b_0\big(\!-\!{\xi\over2}\big)F^{n_2\nu}_{0,b}\big(\!-\!{\xi\over2}\big)\big\rangle_Q\nn\\
	&\equiv\!\int d^4\xi\ \delta(\xi^3\!-\!z) g_{\perp}^{\mu\nu}\big\langle J^{B}_{n_1\mu}\big({\xi\over2}\big)\bar{J}^{B}_{n_2\nu}\big(\!-\!{\xi\over2}\big)\big\rangle_Q\,,\nn
\end{align}
where the auxiliary ``heavy'' quark fields are in the adjoint representation, and
\beq
J^{\mu\nu}_B= F^{\mu\nu}_{0,a}Q_0^a,\ \ \ \ \bar{J}^{\mu\nu}_B= \bar{Q}_0^aF^{\mu\nu}_{0,a}\,.
\eeq

The renormalization of $J^{\mu\nu}_B$ and $\bar{J}^{\mu\nu}_B$ is more involved than the quark case, as they can mix with other composite operators of the same or less dimensions. In DR, BRST symmetry allows $J^{\mu\nu}_B$
to mix with~\cite{Dorn:1986dt,Zhang:2018diq}
\begin{align}
	J^{\mu\nu}_{2B} &\!=\! \left(n_z^\nu F^{\mu n_z}_{0,a} - n_z^\mu F^{\nu n_z}_{0,a} \right) Q_0^a/n_z^2\,,\\
	J^{\mu\nu}_{3B} &\!=\! (-in_z^\mu\! A^\nu_{0,a}\! +\! in_z^\nu A^\mu_{0,a})\big[(in_z\cdot\! D_0\!-\!i\delta m)Q_0\big]^a\!/n_z^2\,.
\end{align}
Their renormalization matrix is given by~\cite{Dorn:1986dt}
\begin{align}\label{eq:oprenormmix}
	\begin{pmatrix}
		J_{B}^{\mu\nu} \\  J_{2B}^{\mu\nu} \\ J_{3B}^{\mu\nu}
	\end{pmatrix}
	&=
	\begin{pmatrix}
		Z_{11} & Z_{12} & Z_{13}\\ 0 & Z_{22} & Z_{23} \\ 0 & 0 & Z_{33}
	\end{pmatrix}
	\begin{pmatrix}
		J_{R}^{\mu\nu}\\  J_{2R}^{\mu\nu} \\ J_{3R}^{\mu\nu}
	\end{pmatrix}
	,
\end{align}
where $J_{2B}^{\mu\nu}$ is gauge invariant while $J^{\mu\nu}_{3B}$ is gauge dependent and proportional to the equation of motion (EOM) for the auxiliary field. The Green's functions of the EOM operator will result in a $\delta$-function,
\begin{align}
	(in_z\cdot D_0(\xi)-i\delta m)\langle Q_0(\xi)\bar{Q}_0(0)\rangle_Q = \delta^{(4)}(\xi)\,,
\end{align}
which only contributes a contact term $\delta(z)$ after integrating over the auxiliary fields. As long as $z\neq0$, such mixing vanishes in all Green's functions of $O^B_g(z)$, so we can ignore the mixing between $J^{\mu\nu}_{B}$ and $J^{\mu\nu}_{3B}$ in the renormalization of $O^B_g(z)$. At $z=0$, $O^B_g(z)$ becomes a local operator and is known to mix with BRST-exact and EOM operators~\cite{Collins:1994ee}, whose renormalization can be performed in the standard way.

Note that when contracted with $n_z$,
\begin{align}
	J_{2B}^{n_z\mu}&\!=\! J_{B}^{n_z\mu} = F^{n_z\mu}_{0,a}Q_0^a\,,\\
	J_{3B}^{n_z\mu}&\!=\! i\big(\!-\!A^{\mu,a}_{0} + {n_z^\mu\over n_z^2} n_z\cdot A_0^a\big)\big[(in_z\cdot D_0-i\delta m)Q_0\big]_a\,,\nn
\end{align}
the $J_{B}^{n_z\mu}$ only mixes with the EOM operator $J_{3B}^{n_z\mu}$. As has been argued above, we can ignore such mixing for $z\neq0$.
Moreover, this degeneracy also leads to relations among elements in the renormalization matrix~\cite{Dorn:1986dt},
\begin{align}\label{eq:Zrelation}
	Z_{11}+Z_{12}=Z_{22}, \,\,\,  Z_{13}=Z_{23}\,.
\end{align}


When contracted with $n_t$,
\begin{align} \label{eq:gauxopt}
	J_{B}^{n_t\mu} &=F^{n_t\mu}_{0,a}Q_0^a\,,\nn\\
	J_{2B}^{n_t\mu} &=n_z^\mu F_{a,0}^{n_tn_z} Q_0^a/n_z^2\,,\nn\\
	J_{3B}^{n_t\mu}&= i{n_z^\mu\over n_z^2} n_t\cdot A_{0}^a \big[(in_z\cdot D_0-i\delta m)Q_0\big]_a\,.
\end{align}
As one can see, $J_{2B}^{n_t\mu}$ and $J_{3B}^{n_t\mu}$ vanish after contraction with $g_{\perp}^{\mu\nu}$, so $J_{B}^{n_t\mu}$ with transverse Lorentz index $\mu$ is multiplicatively renormalizable.

To summarize, for $z\neq0$ and transverse Lorentz index $\mu$, both $J_{B}^{n_z\mu}$ and $J_{B}^{n_t\mu}$ are multiplicatively renormalizable in coordinate space, thus proving the renormalizability of the gluon Wilson-line operator $O^B_g(z)$,
\begin{align} \label{eq:renproofgluon}
	O^B_g(z)&=Z_{J}Z_{\bar{J}}\!\int\! d^4\xi\ \delta(\xi^3\!-\!z)\ g_{\perp}^{\mu\nu}\big\langle J_{n_1\mu}^{R}\big({\xi\over2}\big)\bar{J}_{n_2\nu}^{R}\big(\!-\!{\xi\over2}\big)\big\rangle_Q\nn\\
	&=\ e^{\delta m|z|} Z_{J}Z_{\bar{J}}\ O^R_g(z)\,,
\end{align}
where
\begin{align}
	J_B^{n_1\mu}&=Z_{J}\ J_R^{n_1\mu}=(Z^g_Q)^{1\over 2}Z_A^{1\over 2}Z^g_V J_R^{n_1\mu}\,,\\
	\bar{J}_B^{n_2\nu}&=Z_{\bar{J}}\ J_R^{n_2\nu}=(Z^g_Q)^{1\over 2}Z_A^{1\over 2}Z^g_{\bar{V}} J_R^{n_2\mu}\,,
\end{align}
with $Z^g_V$ and $Z^g_{\bar{V}}$ being the renormalization constants for the vertex involving one gluon and one ``heavy quark'' fields. The wavefunction renormalization constant for the auxiliary ``heavy quark'', $Z^g_Q$, is different from the quark case as it is in the adjoint representation.

In addition, since $J_{B}^{n_z\mu}$ and $J_{B}^{n_t\mu}$ do not mix with ``heavy-to-light'' quark currents due to the mismatch of quantum numbers, it implies that the nonlocal gluon Wilson-line operator does not mix with the singlet quark one under renormalization.

\vspace{.3cm}
For the polarized gluon quasi-PDF, its definition is the same as \eq{quasigluon}, except that the gluon Wilson-line operator becomes
\begin{align}
	\Delta O^B_g(z) = \epsilon_{\perp,\mu\nu} F^{n_1\mu}_{0,a}(z) W^{ab}(z,0) F^{n_2\nu}_{0,b}(0)\,,
\end{align}
where $\epsilon_{\perp}^{\mu\nu}=\epsilon^{03\mu\nu}$. Since $\epsilon_{\perp}^{\mu\nu}$ only contracts with the transverse Lorentz indices, one can use the same proof for $O^B_g(z)$ to show that $\Delta O^B_g(z)$ is also multiplicatively renormalizable and does not mix with singlet quark case~\cite{Zhang:2018diq}.

Finally, one can also prove that \eq{renproofgluon} is valid under lattice regularization with $\delta m$ being linearly divergent~\cite{Zhang:2018diq}. This completes the proof of renormalizability of the gluon Wilson-line operators.

\paragraph*{One-loop renormalization.}
Now we demonstrate the above result by an explicit one-loop example. For the nonlocal Wilson-line operators to be multiplicatively renormalizable, it is important that all linear divergences associated with diagrams other than the Wilson line self-energy cancel out among themselves. To see this, a gauge symmetry preserving regularization is crucial. We use DR and keep poles around $d=3$ to identify the linear divergences~\cite{Zhang:2018diq,Wang:2019tgg}.

The one-loop vertex correction to the ``heavy-to-gluon'' current is shown in \fig{ogmixing}. Each diagram contributes
\begin{align}\label{eq:pole_auxiliary}
	I_a^{\rho\nu}&= \frac{\alpha_s C_A}{\pi}\left[{1\over 4-d}{3\over 4} F_a^{\rho\nu}Q_a + {\rm finite\ terms}\right]\,,\nn\\
	I_b^{\rho\nu}&= \frac{\alpha_s C_A}{\pi}\Big[\frac{1}{d-4}(A_a^\nu n_z^\rho-A_a^\rho n_z^\nu) n_z\cdot \partial Q_a/n_z^2 \nn\\
	&+\frac{\pi\mu}{d-3}\big(n_z^\rho A_a^\nu-n_z^\nu A_a^\rho\big)Q_a+{\rm finite\ terms}\Big]\,,\nonumber\\
	I_c^{\rho\nu}&=\frac{\alpha_s C_A}{\pi}\Big\{\frac{1}{d-4}\Big[\frac{1}{2}\big(F_a^{\rho n_z}n_z^\nu -F_a^{\nu n_z}n_z^\rho \big) Q_a/n_z^2 \nn\\
	&+\frac{1}{4}F_a^{\rho\nu}Q_a+\frac 1 2 (A_a^\rho n_z^\nu-A_a^\nu n_z^\rho) n_z\cdot \partial Q_a/n_z^2\Big]\nonumber\\
	&-\frac{\pi\mu}{d-3}\big(n_z^\rho A_a^\nu-n_z^\nu A_a^\rho\big) Q_a+{\rm finite\ terms}\Big\}\,.
\end{align}

\begin{figure}[htb]	
	\centering
	\begin{subfigure}{.3\linewidth}
		\centering
		\includegraphics[width=\linewidth]{images/pdf-Qg_fig07a.pdf}
		\caption{}
		\label{fig:Qgv1}
	\end{subfigure}
	\begin{subfigure}{.3\linewidth}
		\centering
		\includegraphics[width=\linewidth]{images/pdf-Qg_fig07b.pdf}
		\caption{}
		\label{fig:Qgv2}
	\end{subfigure}	
	\begin{subfigure}{.3\linewidth}
		\centering
		\includegraphics[width=\linewidth]{images/pdf-Qg_fig07c.pdf}
		\caption{}
		\label{fig:Qgv3}
	\end{subfigure}
	\caption{One-loop vertex corrections to the ``heavy-to-gluon'' current.}
	\label{fig:ogmixing}
\end{figure}

Both \fig{Qgv2} and \fig{Qgv3} include a linear divergence that is evident as the $\mu/(d-3)$ term, but they cancel among themselves. This guarantees that the overall UV divergence in the vertex correction is logarithmic, thus the renormalization of the ``heavy-to-gluon'' current is multiplicative. Combining the one-loop results in \eq{pole_auxiliary} and wavefunction renormalizations, we have
\begin{align}
	Z_{11} &= 1+ \frac{\alpha_s C_A}{4\pi}{1\over \epsilon_{\mbox{\tiny UV}}}\,,\qquad\qquad Z_{12} = 1- \frac{\alpha_s C_A}{4\pi}{1\over \epsilon_{\mbox{\tiny UV}}}\,,\nn\\
	Z_{13} &=Z_{23} = 1- \frac{\alpha_s C_A}{4\pi}{1\over \epsilon_{\mbox{\tiny UV}}}\,,\quad Z_{22} = 0\,,
\end{align}
where $C_A=N_c=3$ for QCD.
If we ignore the mixing to the EOM operator,
\begin{align}
	Z^{J^{n_z\nu}}_V &= Z^{J^{\nu n_z}}_V = 0 \,,\nn\\
	Z^{J^{n_ti}}_V &= Z^{J^{in_t}}_V = Z^{J^{ij}}_V = Z^{J^{ji}}_V = 1+ \frac{\alpha_s C_A}{4\pi}  {1\over \epsilon_{\mbox{\tiny UV}}}\,,
\end{align}
where $i,j=1,2$. As a result, the one-loop current renormalization constant is
\begin{align}
	Z_{J^{n_z\nu}} &= Z_{J^{\nu n_z}} =1 +{\alpha_s\over 4\pi}\left({1\over 6}C_A - {4\over 3}n_fT_F\right) {1\over\epsilon_{\mbox{\tiny UV}}} \,,\nn\\
	Z_{J^{n_ti}} &= Z_{J^{in_t}} = Z_{J^{ij}} =Z_{J^{ji}} \nn\\
	&=1 + {\alpha_s\over 4\pi}\left({7\over 6}C_A - {4\over 3}n_fT_F\right) {1\over\epsilon_{\mbox{\tiny UV}}}\,,
\end{align}
where $T_F=1/2$, and $n_f$ is the number of active quark flavors. The two-loop results can be found in~\cite{Braun:2020ymy}.

As one can see, the anomalous dimension of the ``heavy-to-gluon'' current is the same for $\mu,\nu=0,1,2$, which is due to $SO(2,1)$ (or  $SO(3)$ in Euclidean space) symmetry around the $z$-axis.

%------------------------------------------------------------------------------
\subsection{Factorization of Quasi-PDFs}
\label{sec:factorization}
%------------------------------------------------------------------------------


The key to LaMET applications for collinear parton physics
is the factorization formula that relates the quasi-PDFs
to light-cone PDFs~\cite{Ji:2013dva}. Here we use the perturbative properties
of the matching coefficients to write the factorization form
in the $\MS$ scheme in a way consistent with a direct EFT calculations
of PDFs at any given $x$~\cite{Izubuchi:2018srq,Wang:2019tgg}
\begin{align}
	q_i(x,  \mu)\! &=\! \int_{-\infty}^{\infty} \frac{dy}{|y|} \ \bigg[\!\sum_j \tilde C_{q_iq_j}\!\left(\frac{x}{y}, \frac{\mu}{yP^z}\right) \tilde{q}_j(y,P^z,\mu) \label{eq:factorization1} \\
	&\qquad + \tilde C_{qg}\left(\frac{x}{y}, \frac{\mu}{yP^z}\right) \tilde{g}(x,P^z,\mu)\bigg] + \cdots
	\ ,\nn\\
	g(x, \mu)\! &=\! \int_{-\infty}^{\infty} \frac{dy}{|y|} \ \bigg[\!\sum_j \tilde C_{gq}\left(\frac{x}{y}, \frac{\mu}{yP^z}\right) \tilde{q}_j(y,P^z,\mu) \label{eq:factorization2} \\
	&\qquad + \tilde C_{gg}\left(\frac{x}{y}, \frac{\mu}{yP^z}\right) \tilde{g}(y,P^z,\mu)\bigg] + \cdots\,,\nn
\end{align}
where $i,j$ runs over quark and anti-quark flavors.
The ``$\cdots$'' term includes mass corrections whose anayltical forms have been derived to all orders of $M^2/(P^z)^2$~\cite{Chen:2016utp}, and higher-twist contributions of order ${\cal O}\big(\Lambda_{\text{QCD}}^2/ (xP^z)^2,\Lambda_{\text{QCD}}^2/ ((1-x)P^z)^2\big)$ (see \eq{fact_0}). All $P^z$-dependence on the right hand side
cancels out, just like a renormalization scale.

As we have explained in the \sec{lamet}, the above factorization is guaranteed on the physics ground because
the difference between quasi-PDFs and light-cone PDFs is the order of limits in $P^z\to \infty$ and $\Lambda_{\rm UV}\to \infty$, and the IR physics in both quantities must be the same. An all-order
factorization proof for the quark quasi-PDF in perturbation theory was first given with a diagrammatical approach~\cite{Ma:2014jla}. The formula has also been derived using the operator product expansion (OPE) of nonlocal Wilson-line operators~\cite{Izubuchi:2018srq,Ma:2017pxb,Wang:2019tgg}.
Here we outline the diagrammatic proof similar to~\cite{Ma:2014jla}, showing
that the collinear divergences of the quasi-PDFs do factorize and
are equal to those of the light-cone PDFs.
Since the collinear divergence is a concept
in perturbation theory, we will show the factorization using a massless external quark state with lightlike momentum $P^\mu=(P^z,0,0,P^z)$.  While the proof is only for perturbative free quark states,
the factorization formulas are widely believed to be true nonperturbatively as well.
We use DR to regulate both UV and collinear divergences and only consider bare quantities, since UV renormalization does not change the leading collinear divergences.

Before the analysis, we should mention that all the soft divergences cancel between the real and virtual contributions to the quasi-PDFs, as discussed in \sec{lamet-rg}, thus we only need to focus on the collinear divergences. To obtain an intuitive understanding of the structure for collinear divergences, we start from the one-loop diagram in \fig{qoneloop1} in the Feynman gauge. The integral reads
\begin{align}
	\int \frac{d^{4-2\epsilon}k}{(2\pi)^{4-2\epsilon}}\frac{{\rm tr}\big[\slashed{P}\slashed{k}\gamma^z\slashed{k}\big]\delta(k^z-yP^z)}{(k^2+i0)^2((P-k)^2+i0)}\,.
\end{align}
The internal quark momentum is $k^\mu = (k^+,k^-,\vec{k}_\perp)$ and the gluon momentum is $P-k$. When $k^-$ and $k_\perp=|\vec{k}_\perp|$ are very small, the internal quark and gluon become collinear to the external quark, i.e. $k^{\mu}\sim (k^+,0,0_\perp)$ and $(P-k)^{\mu}\sim (P^+-k^+,0,0_\perp)$. In this case, the denominator of the quark and gluon propagators, $(k^2)^2$ and $(P-k)^2$, both vanish, which leads to collinear divergence. Conversely, for $k^2=(P-k)^2=0$, $k$ must be collinear to $P$ since the condition requires $k^2=k\cdot P=P^2=0$. For small $k^-$ and $k_\perp$, the $\delta$ function is dominated by the $k^+$ term of $k^z=(k^+-k^-)/\sqrt{2}$ and reduces to $\sqrt{2}\delta(k^+-yP^+)$. This is just the vertex which restricts $k^+=yP^+$ for the light-cone PDF, up to the factor $\sqrt{2}$. Furthermore, for collinear $k$ and $(P-k)$, the spinor trace in the numerator is dominated by the $\gamma^+$ part of $\gamma^z=(\gamma^+-\gamma^-)/\sqrt{2}$, ${\rm tr}\big[\slashed{P}\slashed{k}\gamma^z\slashed{k}\big] \sim {\rm tr}\big[\slashed{P}\slashed{k}\gamma^+\slashed{k}\big]/\sqrt{2}$. Thus in the collinear region $k^{\mu}\sim (k^+,0,0,0)$ the above integral reduces to that for the light-cone PDF:
\begin{align}
	\int_{\rm c} \frac{d^{4-2\epsilon}k}{(2\pi)^{4-2\epsilon}}\frac{{\rm tr}\big[\slashed{P}\slashed{k}\gamma^+\slashed{k}\big]\delta(k^+-yP^+) }{(k^2+i0)^2((P-k)^2+i0)}\,,
\end{align}
where the subscript ``c'' denotes the collinear region.

The above picture naturally arises in a highly boosted hadron state where the quark is approximately onshell.
Therefore, as explained in \sec{lamet-univ},  {\it although the operator contains no light-cone information, the large-momentum external hadron state can still generate collinear divergences equivalent to those in the light-cone PDFs.} By subtracting the full integral for light-cone PDF from that for the quasi-PDF, the logarithmic collinear divergence cancels, and the remaining difference is perturbative and can be absorbed into the matching kernel.

Similarly, for the vertex diagram in Fig.~\ref{fig:qoneloop2a}, the loop integral is proportional to
\begin{align}
	\int \frac{d^{4-2\epsilon}k}{(2\pi)^{4-2\epsilon}}\frac{1}{P^z-k^z}\frac{{\rm tr}\big[ \slashed{P} \gamma^z \slashed{k} \gamma^z\big] \delta(k^z-yP^z)}{(k^2+i0)((P-k)^2+i0)} \,.
\end{align}
The whole integral in the collinear region reduces to
\begin{align}
	\int_{\rm c} \frac{d^{4-2\epsilon}k}{(2\pi)^{4-2\epsilon}}\frac{1}{P^+-k^+}\frac{{\rm tr} \big[\slashed{P} \gamma^+ \slashed{k} \gamma^+\big]\delta(k^+-yP^+)}{(k^2+i0)((P-k)^2+i0)}  \ ,
\end{align}
which is the corresponding integral for the light-cone PDF. One key feature of the diagram is that while the gauge link probes the $z$-component of the gluon field $A^z =\left(A^+-A^-\right)/\sqrt{2}$, only the $A^+$ component (longitudinal polarization) contributes to the leading collinear divergence. While attaching a new collinear gluon to the gauge-link induces a power suppression from the link propagator of ${\cal O}(1/P^z)$, the $A^+$ component of the collinear gluon radiated from fast-moving color charges receives enhancement from Lorentz boost factor $\gamma$ that compensates for the suppression.

The above result can be generalized to all orders. Similar to the one-loop diagrams, in the leading region of collinear divergence there are an arbitrary number of longitudinally polarized $A^+$ gluons, which are emitted dynamically from the fast-moving state instead of being put in by hand using the lightlike gauge link, in contrast to the standard collinear PDF. The existence of the $A^+$ gluons clearly increases
the level of complication in showing the equivalence of collinear divergences between the quasi- and light-cone PDFs. For simplification, from now on we choose to work in the light-cone gauge $A^+=0$ to eliminate all the $A^+$ gluons. Therefore, the vertex diagrams
no longer contribute to the leading collinear divergence, thus making its structure much simpler.

In a general diagram, we decompose the potential leading region of the quasi-PDF into the ladder structure shown in Fig.~\ref{fig:ladder}. The upper two-particle-irreducible (2PI) kernel that contains the nonlocal operator defining the quasi-PDF is $H$. The 2PI kernel in the ladder is $K$. $K$ contains the upper two external quark lines but not the lower ones. The momentum flowing out of the ladders are labeled as $k_1$ to $k_n$ from bottom to top when there are $n$ 2PI kernels. We write $H$ and $K$ as matrices in spinor and momentum space. $H=H_{\alpha'\beta'}(yP^z;k)$ where $k$ denotes the momentum flowing into $H$ and $K=K_{\alpha\beta;\alpha'\beta'}(k,k')$ where $k$, $k'$ are the momenta of the upper and lower external legs, respectively.
Here $\alpha\beta$ and $\alpha'\beta'$ are the spinor indices for the upper and lower two external legs, respectively.
\begin{figure}
	\includegraphics[width=0.8\columnwidth]{images/pdf-ladder_fig08.pdf}
	\caption{\label{fig:ladder} The ladder decomposition of the quasi-PDF (left). The upper 2PI kernel $H$ contains the operator defining the quasi-PDF, and external two legs at the bottom of the diagram is the external large $P^z$ state. The kernels $H$ and $K$ are shown on the right.}
\end{figure}
Following the method in~\cite{Curci:1980uw,Collins:2011zzd}, we find that:
\begin{enumerate}
	\item There are no collinear divergences in the upper part $H$ in the light-cone gauge.
	\item If none of $k_1,\cdots,k_n$ is collinear, there will be no leading collinear divergence. More generally, for the $i$'th 2PI kernel, if either of $k_{i-1}$ and $k_{i}$ is not collinear, then the sub-integrals inside the kernel are finite and it does not contribute to leading collinear divergence.
	\item If $k_i$ is not collinear, then there are no collinear divergences for the upper part of the diagram above the $i$'th ladder.
\end{enumerate}
Therefore, the collinear divergences are generated in the momentum regions $R_i$ in which $k_1$ to $k_i$ are collinear while $k_{i+1}$ to $k_n$ are not. We can construct counter terms that subtract out the collinear divergences in each of the regions $R_i$.
For this we keep only the $+$ component of $k_i$ in the convergent upper part $HK^{n-i}$ as in the one-loop example, namely $k_i\rightarrow (k^+_i,0,0_\perp)$ in the upper part. This will clearly leave the collinear divergence unchanged.  Also notice that $[HK^{n-i}]_{\alpha\beta}=H_{\alpha'\beta'}K^{n-i}_{\alpha'\beta';\alpha\beta}$ should be understood as a 4 $\times$ 4 Dirac matrix with indices $\alpha\beta$, while the lower part is $[K^i \slashed P]_{\alpha\beta}=K^i_{\alpha\beta;\alpha'\beta'}\slashed P_{\alpha'\beta'}$. In the leading region of collinear divergence, $HK^{n-i}$ and $K^i\slashed P$ are proportional to $\gamma^+$ and $\gamma^-$ respectively.
Therefore, to obtain the leading collinear divergence, we can disentangle the spinor traces for the upper and lower parts by contracting them with $\gamma^-/2$ and $\gamma^+/2$ separately. The only communication between them is the $k^+$ integration. The collinear divergence is contained in the lower part
\begin{align}
	q^{i}(x,\epsilon_{\mbox{\tiny IR}})\!=\!\!\int \frac{dk^-\!d^{d-2}\!k_\perp}{2(2\pi)^{d}}{\rm tr} \big[\gamma^+\! K^i(xP^+\!,k^-\!,k_\perp;P)\slashed{P}\big],
\end{align}
where $d=4-2\epsilon$, $k^+=xP^+$, and the subtraction for the region $R_i$ can be written effectively as a convolution
\begin{align}
	\int \frac{dx}{x} \hat C^{n-i}(y,x,P^z)q^i(x,\epsilon_{\mbox{\tiny IR}}) \ ,
\end{align}
where
\begin{align}
	\hat C^{n-i}(y,x,\!P^z)=\frac{1}{2}{\rm tr}\left[HK^{n-i}(yP^z;xP^+,\!0,0_\perp)(xP^+)\gamma^-\right]
\end{align}
is the naive matching kernel. Here the $y$ dependence comes from the operator in $H$. However, the naive matching kernel still suffers from collinear sub-divergences that need to be subtracted. This can be achieved using the subtracted matching kernels $C^{n-i}(y,x)$ defined recursively in a way similar to the BPHZ relation for UV renormalization~\cite{Collins:2011zzd}. Summing over $n$ and $i$, the recursive relation leads to
\begin{align}\label{eq:subtraction}
	\tilde q(y,P^z,\epsilon_{\mbox{\tiny IR}})&=\sum_{n=0}^{\infty}\sum_{i=0}^n \int \frac{dx}{x} C^{n-i}(y,x,P^z)q^i(x,\epsilon_{\mbox{\tiny IR}})\nn\\
	&=\int \frac{dx}{x} C(y,x,P^z)q(x,\epsilon_{\mbox{\tiny IR}}) \ ,
\end{align}
where $\tilde q(y,P^z,\epsilon_{\mbox{\tiny IR}})$ is the quasi-PDF, $C(y,x,P^z)=\sum_{n=0}^{\infty} C^{n}(y,x,P^z)$ is the all-order matching kernel and $q(x,\epsilon_{\mbox{\tiny IR}})=\sum_{i=0} q^i(x,\epsilon_{\mbox{\tiny IR}})$. Based on the definition of $q^i(x,\epsilon_{\mbox{\tiny IR}})$, it is clear that $q^i$ equals the light-cone PDF with $i$ 2PI kernels and $q$ is the full light-cone PDF with natural support $0<x<1$. {\it The light-cone PDF $q(x)$ is independent of the operator defining the quasi-PDF}, as it is only sensitive to the explicit form of the collinear divergence. The r.h.s. of \eq{subtraction} contains all the collinear divergences from the quasi-PDF $\tilde q$. Thus the matching relation for bare quantities is established. A similar matching can be written down for the renormalized quantities, where the renormalization only affects the matching kernel $C(y,x,P^z)$. Note that the explicit solution for $C^{n-i}(y,x,P^z)$, which leads to \eq{fact_0}, can be given based on a subtraction operator defined similar to that in~\cite{Collins:2011zzd}. Besides, \eq{subtraction} can be inverted order by order in $\alpha_s$, thus proving \eq{fact_0}, which can also be generalized to \eqs{factorization1}{factorization2}.

Now we present the matching coefficient in the $\MS$ scheme at one-loop order.
The one-loop expansion of the $\MS$ quasi- and light-cone PDFs in a free massless quark state with momentum $p^\mu=(p^z,0,0,p^z)$ are
\begin{align}
	\tilde{q}(y,\mu/p^z,\epsilon_{\mbox{\tiny IR}}) &= \tilde{q}^{(0)}(y) + {\alpha_sC_F\over 2\pi}\tilde{q}^{(1)}(y,\mu/p^z,\epsilon_{\mbox{\tiny IR}})\,,\\
	q(x,\epsilon_{\mbox{\tiny IR}}) &= q^{(0)}(x) + {\alpha_sC_F\over 2\pi}q^{(1)}(x,\epsilon_{\mbox{\tiny IR}})\,.
\end{align}
At tree level, $\tilde{q}^{(0)}(y) = q^{(0)}(y) = \delta(1-y)$.
At one loop, the $\MS$ quasi-PDF and its counterterm are~\cite{Izubuchi:2018srq}
\begin{align}\label{eq:qPDFren}
	&\tilde{q}^{(1)}(y,\mu/p^z,\epsilon_{\mbox{\tiny IR}}) \nn\\
	=&\left\{
	\begin{array}{ll}
		\Big({1+y^2\over 1-y}\ln {y\over y-1} + 1 + {3\over 2y}\Big)^{[1,\infty]}_{+(1)}- {3\over 2y} & y>1\\
		\Big({1+y^2\over 1-y}\big[- {1\over \epsilon_{\mbox{\tiny IR}}} - \ln{\mu^2\over 4(p^z)^2} + \ln\big(y(1-y)\big)\big] &\\
		\  - {y(1+y)\over 1-y}+ 2\sigma (1-y)\Big)^{[0,1]}_{+(1)}  &0<y<1\\
		\Big(-{1+y^2\over 1-y}\ln {-y\over 1-y} - 1 + {3\over 2(1-y)}\Big)^{[-\infty,0]}_{+(1)}&\\
		\ - {3\over 2(1-y)} & y<0
	\end{array}\right.\nn\\
	&+ \delta(1-y)\left[{3\over2}\ln{\mu^2\over 4(p^z)^2} + {5 + 2\sigma\over2}\right]\,,
\end{align}
\begin{align}
	\delta \tilde{q}^{(1)}(y,\mu/p^z,\epsilon_{\mbox{\tiny UV}})
	= {3\over 2\epsilon_{\mbox{\tiny UV}}}
	\delta(1-y)\,,
\end{align}
where $\epsilon_{\mbox{\tiny IR}}$ regulates the collinear divergence, $\sigma=0$ for $\Gamma=\gamma^t$ and 1 for $\Gamma=\gamma^z$. The plus function at $y=y_0$ with support in a given domain $D$ is defined as
\begin{align}
	\int_D\! dy \big[ g(y)\big]_{+(y_0)}^D h(y)\! =\! \int_D\! dy\ g(y) \left[ h(y)\! -\! h(y_0)\right]
\end{align}
with arbitrary $g(y)$ and $h(y)$. Note that the $\MS$ renormalization of the quasi-PDF actually requires a subtle treatment of vector current conservation~\cite{Izubuchi:2018srq}. We only present results in the form that is sufficient for our discussion, which differs slightly from that in~\cite{Izubuchi:2018srq} by the $\delta$-functions at $y=\pm\infty$ and from the treatment in~\cite{Alexandrou:2019lfo}.

On the other hand,
\beq\label{eq:lcpdf}
q^{(1)}(x,\epsilon_{\mbox{\tiny IR}}) ={\alpha_sC_F\over 2\pi} {(-1)\over \epsilon_{\mbox{\tiny IR}}} \left({1+x^2\over 1-x}\right)^{[0,1]}_{+(1)}\,,
\eeq
which is limited to the physical region as expected.

By comparing the quasi- and light-cone PDFs in \eqs{qPDFren}{lcpdf}, we find that both of them have the same collinear divergence,  or in other words, they share the same IR physics, thus validating the factorization formula at one-loop order. Setting $p^z=xP^z$ and plugging the one-loop results into \eq{fact_0}, we extract the matching coefficient for the hadron matrix element which only depends on the perturbative scales $\mu$ and $P^z$,
\begin{align}\label{eq:msbarmatching}
	C^{\MS}\left(y, {\mu\over xP^z}\right) &= \delta\left(1-y\right) + {\alpha_sC_F\over 2\pi}\left[\tilde q^{(1)}\left(y,{\mu\over xP^z}, \epsilon_{\mbox{\tiny IR}}\right)\right.\nn\\
	&\qquad\left.-q^{(1)}(y,\epsilon_{\mbox{\tiny IR}})\right]\,.
\end{align}
The complete one-loop matching coefficients in \eq{fact_0} in the transverse-momentum cutoff and $\MS$ schemes can be found in \cite{Wang:2017qyg,Wang:2017eel,Wang:2019tgg}. The two-loop results were obtained recently in~\cite{Chen:2020arf,Chen:2020iqi,Chen:2020ody,Li:2020xml}.

%------------------------------------------------------------------------------
\subsection{Coordinate-Space Factorization of Bilinear Operators}
\label{sec:renorm-coord}
%------------------------------------------------------------------------------
Although the LaMET application to PDFs concerns the expansion of momentum densities
in the $P^z\to\infty$ limit, lattice QCD calculations actually start from
computing coordinate-space correlations, for example,
\beq
\tilde h(z, P^z) = {1\over N_\Gamma}\langle P^z| O_\Gamma(z) |P^z\rangle\,,
\eeq
at all $z$ and do Fourier transform with respect
to $\lambda=zP^z$ at a fixed $P^z$. Here the normalization factor $N_\Gamma=2P^z$ for $\Gamma=\gamma^z$ and $N_\Gamma=2P^t$ for $\Gamma=\gamma^t$.
The $\tilde h(z,P^z)$ is a function of two independent
variables $z$ and $P^z$, and in LaMET analysis the relevant
combinations are quasi-LF distance $\lambda$ (see \fig{Boost}) and $P^z$, hence
$\tilde h(\lambda, P^z)$ will be called quasi-LF correlation, which is distinguished from the LF correlation $h(\lambda,\mu)$ below.

The coordinate-space factorization approach in \cite{Braun:1994jq} has been
suggested as an alternative way to extract the PDFs from $\tilde h(z, P^z)$~\cite{Radyushkin:2017cyf,Orginos:2017kos,Radyushkin:2019mye}, which is closely related to the OPE.
Instead of working with variables $\lambda$ and $P^z$,
one may consider $\tilde h$ as a function of $\lambda$ and $z^2$, i.e.,
$\tilde h(\lambda, z^2)$. The Fourier transform
of $\tilde h(\lambda, z^2)$ with respect to $\lambda$ is no longer the momentum
distribution of the proton at a fixed momentum. Instead, it is called a pseudo-distribution~\cite{Radyushkin:2017cyf}.
At small $|z|\ll \Lambda^{-1}_{\rm QCD}$, $\tilde h(\lambda, z^2)$ can be factorized into the light-cone correlation~\cite{Radyushkin:2017lvu,Izubuchi:2018srq},
\begin{align}
	\label{eq:io-Q-fact}
	\tilde{h}(\lambda, z^2\mu^2)
	=\int_{-1}^1 d\alpha\ \mathcal{C}(\alpha, z^2\mu^2)\ h(\alpha\lambda,\mu) +...\,,
\end{align}
where $...$ are the power corrections in $z^2\Lambda^2_{\rm QCD}$,
and the matching coefficient $\mathcal{C}$
is related to $C$ in Eq.~(\ref{eq:fact_0}) by
\begin{align} \label{eq:ftc-C}
	C\left(\eta,{\mu\over xP^z}\right)
	\!= \!\int {d\lambda\over 2\pi}\ e^{i\eta\lambda}\int_{-1}^1\!\! d\alpha\: e^{-i\lambda\alpha}\: \mathcal{C}\left(\alpha,{\mu^2 \lambda^2\over (xP^z)^2}\right)\,.
\end{align}

To illustrate the connection between the above factorization in \eq{io-Q-fact} and OPE,
let us take the non-singlet quark case as an example~\cite{Izubuchi:2018srq,Wang:2019tgg}. In the $\MS$ scheme, the renormalized $O_{\gamma^{\mu_0}}(z,\mu)$ can be expanded in terms of local gauge-invariant twist-2 operators as $z^2\to 0$,
\begin{align}\label{eq:ope-tq}
	O_{\gamma^{\mu_0}}(z,\mu) &=\sum_{n=0}^\infty \Big[C_n ({\mu}^2 z^2)\frac{(iz)^n}{n!} (n_z)_{\mu_1}\cdots (n_z)_{\mu_n} \nn\\
	&\qquad \times O^{\mu_0\mu_1\cdots\mu_n}(\mu) +  \text{higher-twist}\Big]\,,
\end{align}
where $\mu_0\!=\!0$ or $3$, $C_n\!=\!1+{\cal O}(\alpha_s)$ is the Wilson coefficient, and $O^{\mu_0\mu_1\cdots\mu_n}(\mu)$
is the twist-two operator in \eq{t2opq}.

Using the hadron matrix elements in \eq{t2mom} and their relation to the light-cone PDF
in \eq{t2sum},
we write down the small-$|z|$ expansion of the hadron matrix element of $O_{\gamma^{\mu_0}}(z,\mu) $~\cite{Izubuchi:2018srq},
\begin{align}\label{eq:ope}
	\tilde{h}(\lambda,z^2\mu^2)&=\langle P| O_{\gamma^{\mu_0}}(z,\mu) |P\rangle/(2P^{\mu_0}) \nn\\
	&= \sum_{n=0}^\infty\ C_n(z^2\mu^2){(-i\lambda)^n\over n!}\left[1+{\cal O}\big({M^2\over (P^z)^2}\big)\right]\nn\\
	&\quad\times \int_{-1}^1 dx\ x^n q(x,\mu) + {\cal O}\left(z^2\Lambda_{\text{QCD}}^2\right)\,,
\end{align}
where the ${\cal O}\left(M^2/(P^z)^2\right)$ term comes from the kinematic trace contribution and
the ${\cal O}\left(z^2\Lambda_{\text{QCD}}^2\right)$ term from higher-twist.
The Wilson coefficients $C_n(z^2\mu^2)$ have been calculated at one-loop~\cite{Izubuchi:2018srq} and two-loop~\cite{Li:2020xml} orders.
Comparing the above equation with Eq.~(\ref{eq:io-Q-fact}),
we identify
\begin{align}  \label{eq:fac-C}
	\mathcal{C} (\alpha, {\mu}^2 z^2) \equiv \int\! \frac{d \lambda}{2 \pi} \,
	e^{i  \lambda \alpha}  \sum_n C_n ({\mu}^2 z^2)  \frac{(-i \lambda)^n}{n!}
	\,.
\end{align}
Since $z^2$ is fixed in $\mathcal{C} (\alpha, {\mu}^2 z^2)$, the integration in \eq{fac-C} is actually over $P^z$ from $-\infty$ to $+\infty$.
$\mathcal{C}(\alpha,z^2\mu^2)$ has support $-1 \le \alpha \le 1$,
and its one-loop result is
\begin{align} \label{eq:ps-c}
	&{\cal C}(\alpha,z^2\mu^2)\\
	=& \left[ 1+ {\alpha_sC_F\over 2\pi}\left({3\over2}\ln{z^2\mu^2e^{2\gamma_E}\over 4}+{3\over2}\right)\right]\delta(1-\alpha)\nn
	\\
	&+ {\alpha_sC_F\over 2\pi}\left\{\left({1+\alpha^2\over 1-\alpha}\right)^{[0,1]}_{+(1)}\left[-\ln{z^2\mu^2e^{2\gamma_E}\over 4}-1\right]\right.\nn\\
	&\left. - \left(4\ln(1-\alpha)\over 1-\alpha\right)^{[0,1]}_{+(1)}\!  +\!2(1+\sigma)(1-\alpha)\right\}\!\theta(\alpha)\theta(1-\alpha)
	\,,\nn
\end{align}
which was also calculated and further studied in~\cite{Ji:2017rah,Radyushkin:2017lvu,Radyushkin:2018cvn,Zhang:2018ggy,Li:2020xml}.
One can check that the above result is indeed related to one-loop momentum-space
matching by Eq.~(\ref{eq:ftc-C}).
Since we are interested in the relation between the
quasi-LF correlation with the matrix element of the light-ray operator
$O_{\gamma^+}(\lambda n)$, Eq.~(\ref{eq:io-Q-fact})
can also be obtained by using the light-ray operator
expansion in~\cite{Balitsky:1987bk,Braun:1994jq,Braun:2007wv}.

Using OPE or short-distance expansion, the exact factorization formula for the gluon and singlet quark quasi-PDFs, which includes their mixings, has also been derived in coordinate space and studied at one-loop order~\cite{Wang:2019tgg,Balitsky:2019krf}.

It is easy to see that the limits $P^z\to\infty$ in LaMET expansion and $z\to 0$
in coordinate-space factorization, keeping finite $\lambda=zP^z$, are equivalent.  However, in practical
lattice QCD calculations, one is limited by the largest
momentum $P^z_{\rm max}$ in a specific setup, and the two approaches start to differ.

In LaMET systematic approximation, one should calculate $\tilde h (z, P^z_{\rm max})$ with
all possible $z$ or $\lambda$, but in practice the largest $\lambda_{\rm max} = z_{\rm max}P^z_{\rm max}$
is limited by the lattice volume as well as data quality at large $z$. Due to QCD confinement, $\tilde h (z, P^z_{\rm max})$ has a correlation length $\sim 1/\Lambda_{\rm QCD}$, leading to an exponential decay
at large $z$~\cite{Ji:2020brr}. Therefore, if $z_{\rm max}$ is sufficiently large (e.g., the proton size $\sim1$ fm) for $\tilde h (z, P^z_{\rm max})$ to fall to almost zero, then the truncated Fourier transform of $\tilde h (z, P^z_{\rm max})$
should converge quickly, and the truncation effects mainly affect results at small $x \lesssim 1/\lambda_{\rm max}$. If $\tilde h (z, P^z_{\rm max})$ exhibits exponential decay but still has a nonzero value at $z_{\rm max}$, then one can perform a physically motivated extrapolation beyond $z_{\rm max}$~\cite{Ji:2020brr} to do the Fourier transform, which removes the unphysical oscillation from truncation and only affects the small-$x$ region. In the momentum space, one can use LaMET expansion to calculate the PDF point by point in $x$ with systematic error controlled by $\Lambda_{\rm QCD}^2/(xP^z_{\rm max})^2$ and $\Lambda_{\rm QCD}^2/((1-x)P^z_{\rm max})^2$, which gives the prediction for a cetain region of $x$, $[x_{\rm min}, x_{\rm max}]$, with a target error.

In coordinate-space factorization, one expands $\tilde h(\lambda, z^2)$ in $z^2\Lambda^2_{\rm QCD}$. For the factorization formula to be valid,  $z$ must remain in the perturbative region. For example, an estimate in~\cite{Ji:2020brr} gives $z_{\rm max}\sim 0.3$--$0.4$ fm. Although there have been observations that forming ratios of $\tilde h (\lambda, z^2)$ may cancel the higher-twist contributions at $z>0.4$ fm~\cite{Orginos:2017kos}, this cancellation needs be quantified
for precision calculations. With a finite range of quasi-LF correlations, the PDFs can be extracted through modelling
the $x$-dependence or more advanced techniques such as Bayesian analysis~\cite{Bringewatt:2020ixn} or neural network~\cite{Karpie:2019eiq,Cichy:2019ebf,DelDebbio:2020rgv}, which is similar to extracting the PDFs from experimental data~\cite{Ma:2017pxb}, although quantifying the
systematic error from fitting can be challenging. The coordinate-space factorization can also provide the extraction of the lowest moments of PDFs~\cite{Karpie:2018zaz,Shugert:2020tgq,Joo:2020spy,Gao:2020ito}, where the main systematic error is controlled by $z^2\Lambda_{\rm QCD}^2$.

So far, there have been very limited studies about the comparison between quasi- and pseudo-PDF analysis
~\cite{Bhat:2020ktg,Alexandrou:2020qtt}. It remains to be seen how systematic errors
in the two strategies are compared and contrasted.

%------------------------------------------------------------------------------
\subsection{Nonperturbative Renormalization and Matching}
\label{sec:renorm-npr}
%------------------------------------------------------------------------------
The multiplicative renormalizability of the nonlocal Wilson-line operators for quasi-PDFs allows a nonperturbative renormalization on the lattice, after which the continuum limit can be taken.
This is an important step in the application of LaMET.
One way of doing so is to perform a mass
subtraction of the Wilson line first~\cite{Musch:2010ka,Ishikawa:2016znu,Chen:2016fxx,Zhang:2017bzy,Green:2017xeu,Green:2020xco}, and then renormalize the remnant UV divergences with lattice perturbation theory or other nonperturbative schemes. Another scheme which has gained more popularity in recent years is the regularization-independent momentum subtraction (RI/MOM) scheme~\cite{Constantinou:2017sej,Stewart:2017tvs,Alexandrou:2017huk,Chen:2017mzz,Liu:2018uuj}.
In the coordinate space approach where $|z|\ll \Lambda_{\rm QCD}^{-1}$, the ratios of quasi-LF correlations in
different states~\cite{Radyushkin:2017cyf,Orginos:2017kos,Braun:2018brg,Li:2020xml} have also been proposed as a renormalization scheme. At large $z$, the RI/MOM and ratio schemes introduce extra nonperturbative effects at different levels, which may distort the IR property of the original quasi-LF correlations. Due to the suppression  of long-range contributions by large $P^z$ in the Fourier transform, this nonperturbative contamination mainly affects the end-point region in $x$-space, while the existing LaMET calculations with RI/MOM scheme at moderate $x$, for example in~\cite{Alexandrou:2018eet,Lin:2018qky}, suffers less from such systematics. Nevertheless, the above complication can be avoided by switching to the hybrid scheme~\cite{Ji:2020brr} where one utilizes the advantages of different schemes at short and large distances.
In the following, we discuss the above schemes in order, with a particular focus on the hybrid renormalization scheme.


Before we proceed, it should be noted that the current-current correlators in~\cite{Detmold:2005gg,Braun:2007wv,Ma:2017pxb} do not need or have simple renormalization on the lattice, though it might be more costly to simulate them. Besides, there is another distinct method based on a redefinition of the quasi-PDF with smeared fermion and gauge fields via the gradient flow~\cite{Monahan:2016bvm}. The smeared quasi-PDF is free from UV divergences and remains finite in the continuum limit, which can be perturbatively matched onto the PDF~\cite{Monahan:2017hpu}. Nevertheless, this method awaits to be implemented on the lattice.

\subsubsection{Wilson-line mass-subtraction scheme}
\label{sec:renorm-npr-mass}

Since the mass correction $\delta m$ includes all the linear UV divergences, it is highly favored to nonperturbatively subtract it from the quasi-PDFs. It is well known that the Wilson line renormalization is related to the additive renormalization of the static quark-antiquark potential, i.e., $\delta m$, especially in the context of finite temperature field theory. For a rectangle-shaped Wilson loop of dimension $L\times T$ in the spatial and temporal directions, its vacuum expectation value for large $T$ scales as
\begin{align}
	\lim_{T\to\infty}W(L,T) = c(L) e^{-V(L)T}\,.
\end{align}
The renormalized static potential is
\begin{align}\label{eq:deltam1}
	V^{R}(L) = V(L) + 2\delta m\,,
\end{align}
and $\delta m$ can be fixed by imposing the condition $V^{R}(L_0)=0$ for a particular value of $L_0$. Alternatively, one can also fit $\delta m$ from the famous string potential model,
\begin{align}\label{eq:deltam2}
	V(L) = \sigma L- {\pi \over 12 L} -2\delta m\,.
\end{align}

Apart from using the static potential to determine $\delta m$, it was also proposed to calculate this quantity in the auxiliary ``heavy quark'' field theory with the following condition~\cite{Green:2017xeu},
\begin{align}\label{eq:deltam3}
	\delta m = {d\over dz} \ln{\rm Tr}\big\langle Q(x+zn_z)\bar{Q}(x)\big\rangle_{{\rm QCD}+Q} \Big|_{z=z_0}\,.
\end{align}

Other suggestions have also been made for a nonperturbative calculation of $\delta m$~\cite{Ji:2020brr}. For example, one can investigate the asymptotic large-$z$ behavior of the hadron matrix element or the single quark Green's function, of the vacuum expectation value of $O_{\Gamma}(z,a)$ in a fixed gauge. The $\delta m$ calculated from all these matrix elements will have the following dependence
on the lattice spacing $a$,
\begin{equation}
	\delta m = m_{-1}(a)/a + m_0 \ ,
\end{equation}
where $m_{-1}$(a) is the coefficient of the power divergence which is independent of the specific matrix element, while $m_0~\sim O(\Lambda_{\rm QCD})$ is finite and depends on the external state.
The determination of $m_0$ can be rather nontrivial, and in practical calculations one could adopt a fine-tuning method, such as that for the Wilson-fermion mass, to find the critical value of $m_0$ at which the final result converges fastest in the large $P^z$ limit.

After the Wilson-line mass subtraction, there are still logarithmic UV divergences in $O_\Gamma(z,a)$. One can use lattice perturbation theory to match $\delta m$-subtracted $O_\Gamma(z,a)$ to the $\MS$ scheme~\cite{Ishikawa:2016znu,Xiong:2017jtn,Constantinou:2017sej}, but the convergence still needs to be examined at higher orders. In~\cite{Green:2017xeu,Green:2020xco}, the logarithmic divergences were nonperturbatively renormalized with RI/MOM-like schemes.

The Wilson-line mass-subtraction has been implemented on the lattice in~\cite{Musch:2010ka,Zhang:2017bzy,Green:2017xeu,Chen:2017gck,Alexandrou:2020qtt}.

\subsubsection{RI/MOM scheme}
\label{sec:renorm-npr-rimom}

The RI/MOM scheme has been widely used in lattice QCD for the renormalization of local composite quark operators that are free from power-divergent mixings ~\cite{Martinelli:1994ty}. It is essentially a momentum subtraction scheme in QFT and can be nonperturbatively implemented on the lattice. For an arbitrary composite quark bilinear operator $O^B$ that is multiplicatively renormalized as $O^B=Z_OO^R$, the RI/MOM scheme is defined by imposing the following condition on its off-shell quark matrix element at a subtraction scale $\mu_{\tiny R}$,
\begin{align} \label{eq:rimom}
	Z^{-1}_O \langle p|O^B|p\rangle \Big|_{p^2=-\mu_{\tiny R}^2}= \langle p|O|p\rangle_{\rm tree}\,.
\end{align}
where the subscript ``tree'' means the tree-level matrix element in perturbation theory. If $\mu_{\tiny R} \gg \Lambda_{\rm QCD}$, $Z_O$ defined in \eq{rimom} is in the perturbative region, and we can convert it to the $\MS$ scheme order by order in perturbation theory. In this sense, $Z_O$ is not literally nonperturbative, but an all-order calculable quantity.

Since the nonlocal quark bilinear operator $O_\Gamma(z)$ has been proven to be multiplicatively renormalizable in the coordinate space, one can also renormalize it in the RI/MOM scheme and then match the result to PDF in the $\MS$ scheme~\cite{Constantinou:2017sej,Stewart:2017tvs}.
On the lattice, the off-shell matrix element of an operator is defined from its amputated Green's function, or vertex function, with off-shell quarks. For the nonlocal Wilson-line operator, the latter is
\begin{align} \label{eq:Lambda}
	&\Lambda^{\Gamma}_0(z, a, p) \equiv
	\left[S^{-1}_0(p,a)\right]^\dagger \sum_{x,y}e^{ip\cdot (x-y)}
	\nn\\
	&\ \ \  \times\! \big\langle 0 \big|T\bigl[\psi_0(x,a) O^B_\Gamma(z,a) \bar{\psi}_0(y,a)\bigr]\big|0\big\rangle S^{-1}_0(p,a)\,,
\end{align}
where $S_0(p,a)$ is the bare quark propagator, and the external momentum $p$ is Euclidean on the lattice.
Since Green's functions are not gauge invariant, one needs to fix a gauge (usually Landau gauge $\partial\cdot  A=0$ is chosen), and the gauge dependence is expected to be canceled by the matching or scheme conversion order by order in perturbation theory.

After including the quark wavefunction renormalization $Z_q$, which can be determined independently on the lattice~\cite{Martinelli:1994ty}, \eq{rimom} is revised as
\begin{align} \label{eq:rimomvertex}
	Z_q Z^{-1}_{O_\Gamma} \Lambda^{\Gamma}_0(z,a,p)\Big|_{p=p_{\tiny R}}= \Lambda^{\Gamma}_{\rm tree}(z,a,p) = \Gamma e^{ip_{\tiny R}\cdot z}\,.
\end{align}
Since $O_\Gamma(z,a)$ is not $O(4)$ covariant, one needs to define the RI/MOM scheme with two scales, one is $\mu_{\tiny R}=|p_R|$, and the other $p_{\tiny R}^z$. For convenience we simply denote them as $p=p_{\tiny R}$. To work in the perturbative region and control the lattice discretization effects that are of order ${\cal O}\big(a^2\mu_{\tiny R}^2, a^2(p_{\tiny R}^z)^2\big)$, one must work in the window $\Lambda_{\rm QCD}\ll \mu_{\tiny R}\ll a^{-1}$, $p_{\tiny R}^z \ll a^{-1}$, which is attainable if the lattice spacing is small enough.

Since the quarks are off-shell, also finite mixings with the EOM operators can appear. As a result, \eq{rimomvertex} in general cannot be satisfied as a matrix equation. Instead, one usually needs a projection operator ${\cal P}$ to define the off-shell matrix elements, i.e.
\begin{align}
	\langle p| O^B_\Gamma |p\rangle = {\rm tr}\big[\Lambda^{\Gamma}_0(z,a,p) {\cal P}\big]\,,
\end{align}
so as to calculate the renormalization factor $Z_{O_\Gamma}$.

Then, the bare hadron matrix element $\tilde{h}_B(z,P^z,a)$ can be renormalized in coordinate space as
\begin{align} \label{eq:rimomh}
	\tilde{h}_R(z,P^z, p_{\tiny R}^z,\mu_{\tiny R},a)=Z^{-1}_{O}(z,p_{\tiny R}^z,\mu_{\tiny R},a)\tilde{h}_B(z,P^z,a) \,,
\end{align}
In the continuum limit, the renormalized matrix element is independent of the UV regulator, so we should obtain the same result in DR under RI/MOM scheme, i.e.,
\begin{align}\label{eq:continuum}
	&\tilde{h}_R(z,P^z, p_{\tiny R}^z,\mu_{\tiny R})=\lim_{a\to0}\tilde{h}_R(z,P^z, p_{\tiny R}^z,\mu_{\tiny R},a) \nn\\
	&\qquad = \lim_{\epsilon\to0}Z^{-1}_{O}(z,p_{\tiny R}^z,\mu_{\tiny R},\epsilon)\tilde{h}_B(z,P^z,\epsilon)\,,
\end{align}
which allows us to compute the matching coefficients in continuum perturbation theory. Note that $\delta m$ vanishes in $Z_{O}$ due to the use of DR.

By Fourier transforming the above renormalized matrix element to momentum space, one can then work out the RI/MOM matching coefficient for the quasi-PDFs~\cite{Stewart:2017tvs}. The one-loop matching coefficient for different spin structures has been obtained in \cite{Stewart:2017tvs,Liu:2018uuj,Liu:2018hxv}, and the two-loop result for the unpolarized case can be found in~\cite{Chen:2020ody}. Alternatively, one can also first convert the RI/MOM matrix element to the $\MS$ or modified $\MS$ schemes~\cite{Constantinou:2017sej,Alexandrou:2019lfo}, and then do the Fourier transform and momentum-space matching.

\subsubsection{Ratio scheme}
\label{sec:renorm-npr-ratio}

In the coordinate-space factorization, $|z|\ll \Lambda^{-1}_{\rm QCD}$ must be small, whereas
$P^z$ can be of any value. In this case, the ratio scheme in~\cite{Radyushkin:2017cyf,Orginos:2017kos} can be an effective choice for lattice renormalization.
Consider the ratio
\beq \label{eq:ratio}
\tilde{h}(\lambda,z^2,a)/\tilde{h}(0,z^2,a)\,,
\eeq
where the denominator is a nonperturbative matrix element at $P^z=0$. Since $\tilde{h}(\lambda,z^2,a)$ and $\tilde{h}(0,z^2,a)$ calculated from the same lattice ensemble are correlated with each other, the error in the ratio can be reduced.
Besides, the ratio does not need further renormalization on the lattice, so one can directly take the continuum limit
\begin{align} \label{eq:ratio1}
	\lim_{a\to0}\frac{\tilde{h}(\lambda,z^2,a)}{\tilde{h}(0,z^2,a)} = \frac{\tilde{h}(\lambda,z^2)}{\tilde{h}(0,z^2)}\,,
\end{align}
which has referred to as the ``reduced Ioffe-time pseudo'' distribution in~\cite{Radyushkin:2017cyf,Orginos:2017kos}.
In the $\MS$ scheme, $\tilde{h}(0,z^2\mu^2)$ has a small-$z$ expansion,
\begin{align}
	\tilde{h}(0,z^2\mu^2) = C_0(z^2\mu^2) + {\cal O}(z^2M^2, z^2\Lambda_{\rm QCD}^2)\,,
\end{align}
where the lowest moment of the iso-vector quark PDF $a_0$ is trivially one.
If we ignore all the power corrections, then $\tilde{h}(0,z^2\mu^2)$ is perturbative and can be regarded as a renormalization factor. Therefore, the ratio in \eq{ratio1} still satisfies a similar OPE or factorization formula to \eqs{ope}{io-Q-fact}, except that the matching coefficient must be modified correspondingly~\cite{Radyushkin:2017lvu,Izubuchi:2018srq},
\begin{align}
	{\cal C}^{\rm ratio}(\alpha,z^2\mu^2)= {\cal C}(\alpha,z^2\mu^2) - \delta(1-\alpha)C_0(z^2\mu^2)\,.
\end{align}

In other variants of the ratio scheme, it has also been suggested that one replaces $\tilde{h}(0,z^2,a)$ by the vacuum matrix element of the nonlocal Wilson line operator~\cite{Braun:2018brg,Li:2020xml}, as the UV divergence does not depend on the external state.


\subsubsection{Hybrid scheme}
\label{sec:hybrid}

Since the factorization formula for the quasi-PDF is only proven in the $\MS$ scheme, it is not legitimate to use momentum-space factorization for any other scheme that differ from $\MS$ nonperturbatively. The RI/MOM and ratio schemes fall into this category as the conversion factors that match them to $\MS$ includes logarithms of $z^2$~\cite{Constantinou:2017sej,Izubuchi:2018srq}, which requires running $\alpha_s$ to the IR region when $z\sim \Lambda_{\rm QCD}^{-1}$. In contrast, the Wilson-line mass-subtraction scheme with wavefunction renormalizations is essentially the same as $\MS$, so it will not introduce extra IR effects.

However, the Wilson-line mass-subtraction scheme also has disadvantages. On the lattice, due to discretization effects at $z\sim a$, the lattice scheme cannot reproduce the short-distance $\ln z^2$ behavior of the $\MS$ matrix elements of the nonlocal operator. Such discretization effects, however, are cancelled in the RI/MOM or ratio scheme. To take advantages of both types of schemes, the hybrid scheme was proposed in~\cite{Ji:2020brr} which provides a viable approach to renormalize the quasi-LF correlations at all $z$.


To begin with,
for $|z|\le z_{\rm S}$ where $z_{\rm S}$ is smaller than the distance at which the leading-twist approximation in the OPE becomes unreliable, one renormalizes the quasi-LF correlation as
\begin{align}\label{eq:ratio2}
	\frac{\tilde h(z,a,P^z)}{Z_X(z,a)}\,,
\end{align}
where ``$X$'' can be the RI/MOM or ratio scheme.

For $|z|>z_{\rm S}$, one applies the Wilson-line mass subtraction
\begin{align}
	{\tilde h(z,a,P^z)}e^{-\delta m |z|}Z_{\rm hybrid}\,,
\end{align}
where $Z_{\rm hybrid}$ denotes the wavefunction and vertex renormalizations, which can be nonperturbatively determined by imposing a continuity condition at $z=z_{\rm S}$,
\begin{align}
	Z_{\rm hybrid} e^{-\delta m |z_{\rm S}|}{\tilde h(z,a,P^z)}  = \frac{{\tilde h(z,a,P^z)}}{Z_X(z_{\rm S},a) }\,,
\end{align}
leading to
\begin{align}
	Z_{\rm hybrid}(z_{\rm S},a) = e^{\delta m |z_{\rm S}|}/{Z_X(z_{\rm S},a)  }\,.
\end{align}
In this way, one only has to calculate $\delta m$. Note that the final result should be independent of $z_{\rm S}$, so one should try multiple values and find the optimal one around which the result changes the most slightly.

The perturbative matching for the hybrid renormalized quasi-PDF can be derived accordingly. Taking $Z_X$ being the zero-momentum matrix element in the ratio scheme as an example, the $O(\alpha_s)$ matching has been derived as~\cite{Ji:2020brr}
\begin{align}\label{eq:hybridm2}
	&C_{\rm hybrid}(\xi, \mu^2/(p^z)^2, z_{\rm S}^2\mu^2) = C_{\rm ratio}(\xi, \mu^2/(p^z)^2) \nn\\
	&\quad+{\alpha_sC_F\over 2\pi}{3\over 2} \left[-{1\over |1-\xi|_+}+ \frac{2 \text{Si}((1-\xi) \lambda_{\rm S})}{\pi  (1-\xi)} \right]\,,
\end{align}
where $C_{\rm ratio}$ can be found in~\cite{Izubuchi:2018srq}, $\xi=y/x$, and $\lambda_{\rm S}=z_{\rm S}p^z$ with $p^z=xP^z$ being the parton momentum. The plus function is defined as
\begin{align}
	{1\over |1-\xi|_+} &\equiv \lim_{\beta\to0^+}\left[ {\theta(|1-\xi|-\beta)\over |1-\xi|} + 2\delta(1-\xi)\ln\beta\right] \,.
\end{align}

Due to finite lattice volume and deteriorating signal-to-noise ratios at large $z$, the available lattice data have to be truncated at  $z_{\rm L}$.
As we have discussed in \sec{renorm-coord}, the quasi-LF correlation has a correlation length $\xi_z\sim \Lambda_{\rm QCD}^{-1}$ and exhibits an exponential decay at large $z$ ($\sim 1$ fm). If $z_{\rm L}$ is not sufficiently large and the quasi-LF correlation still has a considerable nonzero value, then a direct Fourier transform truncated at $z_{\rm L}$ will lead to unphysical oscillations and other systematics in the quasi-PDF.

To improve this situation, it is suggested in the hybrid scheme to perform an extrapolation to $z\to\infty$~\cite{Ji:2020brr}. When $P^z$ is not very large and the lattice matrix elements exhibit the exponential behavior near $z_{\rm L}$, one can use the form $\sim e^{-z/\xi_z}$ to do the extrapolation, although some algebraic behavior can be added on top to better reflect the $z$-dependence. If $P^z$ is very large, then the signal-to-noise ratio gets worse, so $z_{\rm L}$ is smaller. In this case, the quasi-LF correlation is yet to show exponential decay and dominated by the leading-twist contributions, so one can use the algebraic decay form to do the extrapolation. Since $\lambda_{\rm L} =z_{\rm L}P^z$ can reach reasonably large values with contemporary computing resources, the extrapolation will only affect very small-$x$ region, for which the LaMET expansion is not well under control after all.

To summarize, the hybrid scheme provides a proper renormalization of the quasi-LF correlations at all $z$, which allows for a controlled calculation of the PDF for $x\in[x_{\rm min}, x_{\rm max}]$ through LaMET expansion in momentum space. Therefore,  we expect it to play a dominant role in the LaMET calculation of PDFs in the future.

\subsection{Total Gluon Helicity $\Delta G$ and Transversity PDF}
\label{sec:deltaG}

Apart from the collinear PDFs, the first application of LaMET is the gluon helicity contribution $\Delta G$ to the proton spin~\cite{Ji:2013fga}. In the naive sum rule for the proton spin~\cite{Jaffe:1989jz}, $\Delta G$ is related to the matrix element of a nonlocal light-cone correlation operator~\cite{Manohar:1990jx},
\begin{align}\label{eq:sginv}
	\Delta G &\!=\! \langle PS\big|i\!\int\! \frac{dxd\lambda}{2\pi xP^+} e^{i\lambda x} \! F^{+\alpha}(0) W(0,\!\lambda n)\tilde F_{\alpha}^{~+} (\lambda n)\big|PS\rangle\,,
\end{align}
which in the light-cone gauge $A^+=0$ reduces to
\begin{align}
	\Delta G &= \langle PS| \big(\vec{E}\times \vec{A}\big)^z|PS\rangle /(2P^+)\,.
\end{align}

Within the LaMET framework, one can start from  a static ``gluon spin'' operator, which is defined as $\vec{E}\times \vec{A}$ fixed in a time-independent gauge which maintains
the transverse polarizations of the gluon field in the IMF limit. For example, the Coulomb gauge $\vec{\nabla}\cdot \vec{A}=0$, axial gauges $A^z=0$ and $A^0=0$ are viable options~\cite{Hatta:2013gta}.

In the Coulomb gauge and $\MS$ scheme, the static ``gluon spin'' $\Delta \widetilde G$ in a massive on-shell quark state at one-loop order is~\cite{Chen:2011gn,Ji:2013fga}
\begin{align}  \label{eq:coulomb}
	&\Delta \widetilde G(P^z, \mu)  (2S^z)=\left.\langle PS| \epsilon^{ij}_\perp F^{i0}A^j|PS\rangle_q \right\arrowvert_{\vec{\nabla}\cdot \vec{A}=0}\\
	&\qquad\qquad = \frac{\alpha_sC_F}{4\pi} \left[ {5\over 3}\ln{\mu^2\over m^2} -\frac{1}{9} + \frac{4}{3}\ln \frac{(2P^z)^2}{m^2}\right](2S^z) \,,\nn
\end{align}
where the subscript $q$ denotes a quark, and $S^\mu$ is the spin vector. The collinear divergence is regulated by the finite quark mass $m$.

If we follow the procedure in~\cite{Weinberg:1966jm} and take $P^z\to\infty$ limit before UV regularization~\cite{Ji:2013fga}, then
\begin{align} \label{eq:coulombimf}
	\Delta \widetilde{G}(\infty, \mu) (2S^z)&= \left. \langle PS| \epsilon^{ij}_\perp F^{i0}A^j|PS\rangle_q\right\arrowvert_{\vec{\nabla}\cdot \vec{A}=0}\nn\\
	&=\frac{\alpha_sC_F}{4\pi}\left(3\ln{\mu^2\over m^2}+7\right) (2S^z)\,,
\end{align}
which is exactly the same as the light-cone gluon helicity $\Delta G(\mu)$~\cite{Hoodbhoy:1998bt}. Therefore,
despite the difference in the UV divergence, the collinear divergences of $\Delta \widetilde G(P^z,\mu)$ and $\Delta G(\mu)$ are exactly the same,
which allows for a perturbative matching between them.

The complete factorization formula that relates $\Delta\widetilde G(P^z,\mu)$ to $\Delta G$ and $\Delta \Sigma$ is
\begin{align} \label{eq:deltaGmatching}
	\Delta\widetilde G(P^z,\mu) &= Z_{gg}(P^z/\mu)\Delta G(\mu) \nn\\
	&\qquad\qquad+ Z_{gq}(P^z/\mu)\Delta\Sigma(\mu) +...\,,
\end{align}
where $...$ are power corrections suppressed by $1/P^z$, and the matching coefficients $Z_{gg}$ and $Z_{gg}$ have been calculated for the Coulomb gauge at one-loop~\cite{Ji:2014lra}.

Besides, one can also calculate the gluon helicity PDF $\Delta g(x)$
according to the factorization formula in \sec{renorm}, and then integrate it over $x$ to obtain $\Delta G$.\\

At leading-twist, apart from the unpolarized
and helicity PDFs that we have discussed before, there is also the transversity
PDF defined as~\cite{Jaffe:1991kp,Jaffe:1991ra}
\begin{equation}
	h_1(x)\! =\! \frac{1}{2P^+}\!\int \frac{d\lambda}{2\pi}
	e^{-i\lambda x} \langle PS_\perp|\overline{\psi}(0)\gamma^+\gamma_\perp\gamma_5\psi(\lambda n)|PS_\perp\rangle \,.
\end{equation}
The $h_1(x)$ simply counts the number of transversely polarized quarks carrying
the momentum fraction $x$ in a transversely polarized proton.
The first moment of this distribution corresponds to the
so-called tensor charge $\delta q$, which is the matrix element of a chiral-odd operator. $h_1(x)$ can be accessed through the transverse-transverse spin asymmetry in Drell-Yan processes~\cite{Ralston:1979ys,Jaffe:1991kp,Jaffe:1991ra} or the Collins single-spin asymmetry in SIDIS where the transversity TMDPDF couples to a chiral-odd TMD fragmentation function~\cite{Collins:1992kk}.
At present, experimental results on the transversity PDF are very
limited~\cite{Barone:2001sp,Kang:2015msa,Lin:2017stx,Radici:2018iag,Cammarota:2020qcw}, especially for the sea quark contributions~\cite{Chang:2014jba}, so this is one scenario where lattice
QCD calculation can make an important difference.

The LaMET calculation of $h_1(x)$ is straightforward as the nonlocal operator has the same renormalization as the unpolarized case, and its one-loop matching has been calculated in the $\MS$ and RI/MOM schemes at one-loop order~\cite{Alexandrou:2018eet,Liu:2018hxv}. First lattice calculations of $h_1(x)$ have been done in~\cite{Chen:2016utp,Alexandrou:2018eet,Liu:2018hxv}, which will be discussed with more details in \sec{lattice}.


%------------------------------------------------------------------------------
